\documentclass[11pt]{article}
\usepackage[margin=1in]{geometry}
\usepackage{hyperref}
\usepackage{listings}
\usepackage{xcolor}
\usepackage{longtable}
\usepackage{booktabs}
\usepackage{enumitem}
\usepackage{amsmath}

\lstset{
    basicstyle=\ttfamily\small,
    backgroundcolor=\color{gray!10},
    frame=single,
    framerule=0.5pt,
    breaklines=true,
    showstringspaces=false,
    columns=flexible
}

\title{FEMM Current Flow File Formats\\[0.5em]
\large Input (\texttt{.fec}) and Output (\texttt{.anc}) Specifications}
\author{D.\ Meeker}
\date{\today}

\begin{document}
\maketitle
\tableofcontents
\newpage

% ======================================================================
\section{Overview}

The FEMM current flow solver (\texttt{csolv}) reads a problem
definition from a \texttt{.fec} file and writes the solution to a
\texttt{.anc} file.  The \texttt{.anc} file begins with a verbatim
copy of the input \texttt{.fec} file, followed by a
\texttt{[Solution]} section containing the finite element mesh and
computed complex voltage values.

Current flow is an AC problem.  The solver computes the complex
voltage phasor $\hat{V}$ at each mesh node, from which the electric
field $E = -\nabla V$ and current density $J = \sigma E + j\omega
\varepsilon E$ are derived.  All solution quantities (voltages,
currents, fields) are complex-valued.

All floating-point values are written with 17 significant digits
(\verb|%.17g| format).  Fields within a line are separated by tabs or
spaces.  Section headers use square brackets (\verb|[Tag]|); property
block delimiters use angle brackets (\verb|<BeginBlock>| /
\verb|<EndBlock>|).  String comparisons during parsing are
case-insensitive.

% ======================================================================
\section{The \texttt{.fec} Input File}

% ----------------------------------------------------------------------
\subsection{Global Problem Parameters}

Each parameter occupies one line in the format
\verb|[Tag] = value|.  All are optional; defaults are applied if
missing.

\begin{longtable}{lllp{5.5cm}}
\toprule
Tag & Type & Default & Description \\
\midrule
\endhead
\verb|[Format]|      & int     & 1          & File format version \\
\verb|[Frequency]|   & double  & 60         & AC frequency (Hz) \\
\verb|[Precision]|   & double  & 1e-8       & Solver convergence tolerance \\
\verb|[MinAngle]|    & double  & 30         & Minimum mesh element angle (degrees) \\
\verb|[DoSmartMesh]| & int     & 1          & Adaptive meshing (0=off, 1=on) \\
\verb|[Depth]|       & double  & 1.0        & Extrusion depth (planar problems) \\
\verb|[LengthUnits]| & string  & millimeters & inches, millimeters, centimeters, meters, mils, microns \\
\verb|[ProblemType]| & string  & planar     & planar or axisymmetric \\
\verb|[Coordinates]| & string  & cartesian  & cartesian or polar \\
\verb|[Comment]|     & string  & (empty)    & Problem description (quoted; \verb|\n| = newline) \\
\verb|[ACSolver]|    & int     & 0          & Always 0 (compatibility field) \\
\bottomrule
\end{longtable}

The \verb|[Frequency]| field distinguishes current flow from
electrostatics.  At $f = 0$, the displacement current term vanishes
and the problem reduces to pure DC conduction.  The solver internally
converts frequency to angular frequency: $\omega = 2\pi f$.

\subsubsection{Axisymmetric External Region}

Written only when \verb|ProblemType| is axisymmetric and both
\verb|extRo| and \verb|extRi| are nonzero:

\begin{longtable}{llp{8cm}}
\toprule
Tag & Type & Description \\
\midrule
\verb|[extZo]| & double & Outer Z coordinate \\
\verb|[extRo]| & double & Outer radius \\
\verb|[extRi]| & double & Inner radius \\
\bottomrule
\end{longtable}

% ----------------------------------------------------------------------
\subsection{Point (Nodal) Properties}

All point property values are complex (separate real and imaginary
parts) to support AC excitation.

\begin{lstlisting}
[PointProps] = <count>
<BeginPoint>
  <PointName> = "name"
  <vpr> = 0        prescribed voltage, real (V)
  <vpi> = 0        prescribed voltage, imag (V)
  <qpr> = 0        point current source, real (A)
  <qpi> = 0        point current source, imag (A)
<EndPoint>
... repeated count times
\end{lstlisting}

% ----------------------------------------------------------------------
\subsection{Boundary Properties}

All boundary values are complex-valued.

\begin{lstlisting}
[BdryProps] = <count>
<BeginBdry>
  <BdryName> = "name"
  <BdryType> = 0
  <vsr> = 0          fixed voltage, real (V)
  <vsi> = 0          fixed voltage, imag (V)
  <qsr> = 0          surface current density, real (A/m)
  <qsi> = 0          surface current density, imag (A/m)
  <c0r> = 0          mixed BC coefficient, real
  <c0i> = 0          mixed BC coefficient, imag
  <c1r> = 0          mixed BC coefficient, real
  <c1i> = 0          mixed BC coefficient, imag
<EndBdry>
\end{lstlisting}

Boundary types:
\begin{itemize}[nosep]
  \item 0 = fixed voltage (Dirichlet): $V = V_s$
  \item 1 = mixed (Robin): $\hat{n} \cdot J = c_0 + c_1 V$
  \item 2 = surface current density (Neumann)
  \item 3 = periodic
  \item 4 = antiperiodic
\end{itemize}

% ----------------------------------------------------------------------
\subsection{Material (Block) Properties}

Current flow materials combine electrical conductivity and dielectric
properties.  The solver constructs an effective complex admittivity
from these parameters:
\[
  K_x = \frac{\sigma_x}{\varepsilon_0} + j\omega\,\varepsilon_{r,x}
        \, e^{-j\arctan(\tan\delta_x)}
\]
and similarly for $K_y$.

\begin{lstlisting}
[BlockProps] = <count>
<BeginBlock>
  <BlockName> = "name"
  <ox> = 0          DC conductivity, x-direction (MS/m)
  <oy> = 0          DC conductivity, y-direction (MS/m)
  <ex> = 1          relative permittivity, x-direction
  <ey> = 1          relative permittivity, y-direction
  <ltx> = 0         loss tangent, x-direction (tan delta)
  <lty> = 0         loss tangent, y-direction (tan delta)
<EndBlock>
\end{lstlisting}

\begin{itemize}[nosep]
  \item Conductivity $\sigma$ is in megasiemens per meter (MS/m).
  \item Loss tangent $\tan\delta$ creates a complex permittivity that
        models dielectric loss.  Zero means lossless dielectric.
  \item At $f = 0$ (DC), the displacement current term vanishes and
        only the conductivity term remains.
  \item At $\sigma = 0$ (perfect insulator), only displacement current
        flows.
\end{itemize}

% ----------------------------------------------------------------------
\subsection{Conductor Properties}

Conductors enforce equipotential constraints or prescribe total
current on electrode surfaces.  All values are complex.

\begin{lstlisting}
[ConductorProps] = <count>
<BeginConductor>
  <ConductorName> = "name"
  <vcr> = 0           prescribed voltage, real (V)
  <vci> = 0           prescribed voltage, imag (V)
  <qcr> = 0           prescribed total current, real (A)
  <qci> = 0           prescribed total current, imag (A)
  <ConductorType> = 0
<EndConductor>
\end{lstlisting}

Conductor types:
\begin{itemize}[nosep]
  \item 0 = prescribed total current ($\hat{q}$ given; solver
        computes $\hat{V}$)
  \item 1 = prescribed voltage ($\hat{V}$ given; solver computes
        $\hat{q}$)
\end{itemize}

% ----------------------------------------------------------------------
\subsection{Geometry}

\subsubsection{Nodes}

\begin{lstlisting}
[NumPoints] = <count>
x    y    propIdx    inGroup    conductorIdx    [lfs]
\end{lstlisting}

\begin{itemize}[nosep]
  \item \texttt{x}, \texttt{y}: coordinates (double)
  \item \texttt{propIdx}: index into point property list (1-based; 0 = none)
  \item \texttt{inGroup}: group number (int)
  \item \texttt{conductorIdx}: index into conductor list (1-based; 0 = none)
  \item \texttt{lfs}: local feature size constraint (double, optional; $-1$ = auto)
\end{itemize}

\subsubsection{Line Segments}

\begin{lstlisting}
[NumSegments] = <count>
n0    n1    maxSideLength    propIdx    hidden    inGroup    conductorIdx
\end{lstlisting}

\begin{itemize}[nosep]
  \item \texttt{n0}, \texttt{n1}: endpoint node indices (0-based)
  \item \texttt{maxSideLength}: mesh size constraint (double; $-1$ = auto)
  \item \texttt{propIdx}: boundary property index (1-based; 0 = none)
  \item \texttt{hidden}: 0 = visible, 1 = suppressed from mesh
  \item \texttt{inGroup}: group number
  \item \texttt{conductorIdx}: conductor index (1-based; 0 = none)
\end{itemize}

\subsubsection{Arc Segments}

\begin{lstlisting}
[NumArcSegments] = <count>
n0  n1  arcAngle  maxSideLength  propIdx  hidden  inGroup  conductorIdx  mySideLength
\end{lstlisting}

\begin{itemize}[nosep]
  \item \texttt{n0}, \texttt{n1}: endpoint node indices (0-based)
  \item \texttt{arcAngle}: arc span (degrees)
  \item \texttt{maxSideLength}: mesh constraint (double)
  \item \texttt{propIdx}: boundary property index (1-based; 0 = none)
  \item \texttt{hidden}: 0 = visible, 1 = suppressed
  \item \texttt{inGroup}: group number
  \item \texttt{conductorIdx}: conductor index (1-based; 0 = none)
  \item \texttt{mySideLength}: computed side length (0 = auto)
\end{itemize}

\subsubsection{Holes}

\begin{lstlisting}
[NumHoles] = <count>
x    y    inGroup
\end{lstlisting}

\subsubsection{Block Labels}

\begin{lstlisting}
[NumBlockLabels] = <count>
x    y    blockIdx    meshSize    inGroup    extFlags
\end{lstlisting}

\begin{itemize}[nosep]
  \item \texttt{x}, \texttt{y}: label position (double)
  \item \texttt{blockIdx}: material index (1-based; 0 = unassigned)
  \item \texttt{meshSize}: element size (double; $\leq 0$ = auto).
        Stored as diameter; converted to area: $A = \pi d^2/4$
  \item \texttt{inGroup}: group number
  \item \texttt{extFlags}: bit 0 = external region, bit 1 = default region
\end{itemize}

% ----------------------------------------------------------------------
\subsection{Differences from Other Physics Formats}

\begin{itemize}
  \item \texttt{[Frequency]} is always present and typically nonzero
        (default 60~Hz), unlike electrostatics (always 0) and heat
        flow (no frequency).
  \item All point properties, boundary values, and conductor
        properties have separate real/imaginary parts (complex AC
        phasors).
  \item Material properties include both conductivity ($\sigma_x$,
        $\sigma_y$) and permittivity ($\varepsilon_x$,
        $\varepsilon_y$) with loss tangent --- unique to current flow.
  \item No nonlinear material curves (no B-H or T-K equivalent).
  \item No transient fields (\texttt{[dT]}, \texttt{[PrevSoln]} are
        heat-flow-specific).
\end{itemize}


% ======================================================================
\section{The \texttt{.anc} Output File}

The \texttt{.anc} file begins with a verbatim copy of the input
\texttt{.fec} file (all sections described in Section~2), followed by
the solver's solution data.

% ----------------------------------------------------------------------
\subsection{Solution Header}

Immediately after the echoed \texttt{.fec} content:

\begin{lstlisting}
[Solution]
\end{lstlisting}

% ----------------------------------------------------------------------
\subsection{Mesh Nodes}

\begin{lstlisting}
<NumNodes>
\end{lstlisting}

followed by \texttt{NumNodes} lines:

\begin{lstlisting}
x    y    V_re    V_im    Q
\end{lstlisting}

\begin{itemize}[nosep]
  \item \texttt{x}, \texttt{y}: node coordinates in model units
  \item \texttt{V\_re}: real part of complex voltage phasor (V)
  \item \texttt{V\_im}: imaginary part of complex voltage phasor (V)
  \item \texttt{Q}: boundary/conductor marker (int):
    \begin{itemize}[nosep]
      \item $-2$ = interior node (no boundary condition)
      \item $-1$ = fixed voltage boundary
      \item $0 \ldots N-1$ = conductor index
    \end{itemize}
\end{itemize}

The voltage phasor represents the steady-state AC solution.  The
instantaneous voltage at any point is $V(t) = |\hat{V}| \cos(\omega t
+ \arg\hat{V})$.

Node indices are 0-based throughout the solution section.

% ----------------------------------------------------------------------
\subsection{Mesh Elements}

\begin{lstlisting}
<NumElements>
\end{lstlisting}

followed by \texttt{NumElements} lines:

\begin{lstlisting}
p0    p1    p2    labelIdx
\end{lstlisting}

\begin{itemize}[nosep]
  \item \texttt{p0}, \texttt{p1}, \texttt{p2}: vertex node indices (0-based)
  \item \texttt{labelIdx}: block label index (0-based, into the block
        label list from the \texttt{.fec} header)
\end{itemize}

% ----------------------------------------------------------------------
\subsection{Conductor Solution Data}

\begin{lstlisting}
<NumConductors>
\end{lstlisting}

followed by \texttt{NumConductors} lines:

\begin{lstlisting}
V_re    V_im    q_re    q_im
\end{lstlisting}

\begin{itemize}[nosep]
  \item \texttt{V\_re}, \texttt{V\_im}: complex voltage on the
        conductor (V).  For prescribed-voltage conductors, this is
        the prescribed value; for prescribed-current conductors, this
        is the computed value.
  \item \texttt{q\_re}, \texttt{q\_im}: complex total current through
        the conductor (A).  For prescribed-current conductors, this is
        the prescribed value; for prescribed-voltage conductors, this
        is the computed value.
\end{itemize}

% ======================================================================
\section{Index Conventions}

\begin{itemize}
  \item \textbf{Property indices} (materials, boundaries, conductors,
        point properties) are \textbf{1-based} in the file.  An index
        of 0 means ``no property assigned.''  The parser converts to
        0-based internally.

  \item \textbf{Node indices} (in segments, arcs, and the solution
        section) are \textbf{0-based}.

  \item \textbf{Block label indices} in the solution element list are
        \textbf{0-based}, referencing the block label array from the
        \texttt{.fec} header.
\end{itemize}

% ======================================================================
\section{Unit Conversion}

The \verb|[LengthUnits]| tag determines the coordinate system.  All
coordinates in the file are in the specified units.  Conversion to
meters:

\begin{center}
\begin{tabular}{ll}
\toprule
Unit & Meters per unit \\
\midrule
inches      & 0.0254 \\
millimeters & 0.001 \\
centimeters & 0.01 \\
meters      & 1.0 \\
mils        & $2.54 \times 10^{-5}$ \\
microns     & $10^{-6}$ \\
\bottomrule
\end{tabular}
\end{center}

\end{document}
